\section{Normed chain complexes: repairing homotopies from a finite defect}
\label{sec:chain-repair}

The useful categorical output in this section is an explicitly represented
inverse \emph{up to chain homotopy}.  It is obtained from a quantitative
certificate, without first constructing homology groups, selecting quotient
representatives, or deciding membership in an image.  The underlying
perturbation idea is classical; see Crainic~\cite{crainic2004}, especially
Sections~2.2--2.3 and~3.4.  What matters here is its organization into exact
finite identities, rational bounds, a primitive recursive precision schedule,
and a precisely delimited iPRA application.

The existing categorical material already distinguishes fixed presentations,
finite tables, and an external category of specified programs
\cite{repo-categories}.  We retain that distinction.  Its obstruction to
selecting representatives merely from a bound on the number of components
\cite{repo-obstacles} is not removed by the construction below.  Instead,
an explicit contracting operator supplies the missing witness in a stable
analytic situation.

\subsection{Data and meaning of the norm bounds}

Fix a finite band of degrees and a graded rational normed vector space
$C_0=\bigoplus_i(C_0)_i$.  Each summand may be infinite dimensional.
Valid core vectors, addition, rational scalar multiplication, and a rational
norm are given by fixed PR recipes.  Their algebraic and norm laws are
supplied with the arithmetic proofs required by the framework.  Write
$\widehat C$ for the space of represented Cauchy names over this core.
This notation is external shorthand for the represented objects; it does
not quantify in iPRA over all sequences or over all real vectors.

A differential is a fixed linear recipe $D$ of degree $-1$, satisfying
$D^2x=0$ and $\|Dx\|\le B_D\|x\|$ on valid core vectors.  A candidate
contracting homotopy is a fixed linear recipe $h$ of degree $+1$, with
$\|hx\|\le H\|x\|$.  At the ends of the degree band the unavailable
components are zero.  The bounds extend these maps to represented
completions.  Throughout this section, notation such as $\|T\|\le b$
means the supplied assertion
\[
                         \|Tx\|\le b\|x\|
                         \quad\hbox{for every valid }x.
\]
It does not require computing an operator norm as an unbounded supremum.
Only the fixed recipes that occur in the proof are used.

The core-preserving case gives a direct PR construction.  The same finite
proof applies to represented operators that take core points to Cauchy
names when their finite compositions are supplied with evaluation and error
contracts.  In that version the error ledger must pay for those evaluations;
the algebraic identities alone do not supply them.  In particular, arbitrary
bounded operators are not silently included in a uniform evaluation domain.

Define the degree-zero defect and its finite correction by
\begin{equation}
 R=I-Dh-hD,\qquad
 S_N=\sum_{j=0}^{N-1}R^j,\qquad h_N=hS_N,\qquad S_0=0.
 \label{eq:chain-defect}
\end{equation}
All products mean composition, with the rightmost operator applied first.
The candidate need not satisfy $h^2=0$, and $h$ need not commute with $R$.

\begin{lemma}[Exact finite correction]
\label{lem:finite-chain-correction}
For every natural number $N$,
\begin{equation}
 DR=RD,\qquad DS_N=S_ND,\qquad
                Dh_N+h_ND=I-R^N.
 \label{eq:chain-finite-identity}
\end{equation}
\end{lemma}
\begin{proof}
Expanding the first equation and using $D^2=0$ gives
\[
 DR=D-D^2h-DhD=D-DhD
    =D-DhD-hD^2=RD.
\]
Induction on $j$ gives $DR^j=R^jD$: the successor calculation is
$DR^{j+1}=(DR^j)R=R^jDR=R^{j+1}D$.
Adding these equations proves $DS_N=S_ND$.  The two finite telescoping
identities are
\[
 (I-R)S_N=I-R^N=S_N(I-R).
\]
They hold for $N=0$ because both sides are zero.  Adding $R^N$ to $S_N$
adds $R^N-R^{N+1}$ to either product, which proves the successor step.
Finally,
\[
 Dh_N+h_ND
   =DhS_N+hS_ND
   =(Dh+hD)S_N
   =(I-R)S_N=I-R^N.
\]
Every step is an equation of finite terms.  In particular, this proof does
not use convergence, commutativity of the operator algebra, or a side
condition on $h^2$.
\end{proof}

\subsection{A fully explicit correction schedule}

Assume a rational defect certificate
\begin{equation}
      0\le\rho\le 1-2^{-s},\qquad
      \|R\|\le\rho,\qquad H\le 2^c,
      \qquad s,c\in\mathbb N.
 \label{eq:chain-margin}
\end{equation}
The non-strict rational inequality contains a positive margin even when the
actual operator norm is unknown.  A bare assertion $\|R\|<1$ must first
provide such a usable margin before it can serve as a numerical contract.

For $M\ge N$, finite expansion gives
\begin{align}
 \|h_N\|&\le H\sum_{j<N}\rho^j
                  \le\frac{H}{1-\rho},\label{eq:chain-bound}\\
 \|h_M-h_N\|&\le H\sum_{j=N}^{M-1}\rho^j
                  \le\frac{H\rho^N}{1-\rho}.
                  \label{eq:chain-tail}
\end{align}
These are finite geometric-sum estimates.  For example,
$(1-\rho)\sum_{j<N}\rho^j=1-\rho^N\le1$; multiplication by the
positive denominator justifies the displayed rational bound.  Bounds for
products follow from applying the two operator inequalities successively
to an arbitrary vector.  Thus no supremum principle is hidden here.

Here is a schedule that uses no logarithm, exponential approximation,
or unbounded least search.  For $0\le a\le1$, induction on $m$ proves
\begin{equation}
                   (1-a)^m(1+ma)\le1.
 \label{eq:rational-geometric-block}
\end{equation}
Indeed,
$(1-a)(1+(m+1)a)=1+ma-(m+1)a^2\le1+ma$.
Taking $a=2^{-s}$ and $m=2^s$ gives
$\rho^{2^s}\le1/2$, hence
\[
                         \rho^{2^s t}\le2^{-t}.
\]
Consequently
\begin{equation}
                  N(n)=2^s(n+s+c+2)
                  \label{eq:chain-schedule}
\end{equation}
makes the right side of \eqref{eq:chain-tail} at most $2^{-n-2}$.
After denominators are cleared, the induction invariant in
\eqref{eq:rational-geometric-block} is a quantifier-free inequality of
PR integer terms.  This is an explicit use of quantifier-free induction,
rather than induction over a proposition about all completed vectors.

\begin{theorem}[Repair of a represented contracting homotopy]
\label{thm:chain-repair}
Under the stated presentation laws and \eqref{eq:chain-margin}, the recipes
$h_N$ define a represented bounded linear map $\mathcal H$ on
$\widehat C$ satisfying
\begin{equation}
 \|\mathcal H\|\le\frac{H}{1-\rho}\le2^{s+c},\qquad
 D\mathcal H+\mathcal H D=I,
 \label{eq:chain-repaired}
\end{equation}
and
\begin{equation}
          \|\mathcal H-h_N\|
                   \le\frac{H\rho^N}{1-\rho}.
          \label{eq:chain-limit-tail}
\end{equation}
The map and all its requested observations have fixed PR recipes.
\end{theorem}
\begin{proof}
For a fixed core vector $x$, the estimate \eqref{eq:chain-tail}, multiplied
by a supplied bound on $\|x\|$, gives a PR Cauchy modulus for $h_Nx$.
The completion construction therefore represents its limit.
There is also an explicit diagonal formula for a non-core input.
Suppose $x_m$ is the fixed PR name of $x$, with
$\|x-x_m\|\le2^{-m}$ and $\|x\|\le2^b$.  Set
\[
 m=n+s+c+3,\qquad
 N=2^s(n+b+s+c+3),\qquad y_n=h_Nx_m.
\]
The error from the input name is at most
$2^{s+c}2^{-m}=2^{-n-3}$.  The truncation error on $x$ is at most
$2^{s+c+b}\rho^N\le2^{-n-3}$.  Thus $y_n$ is a rational-core
observation of the desired output with error at most $2^{-n-2}$.
The use of $x_m$ is schematic over this fixed name program, not an
evaluation operation for arbitrary PR program codes.

Each $h_N$ is linear.  To prove linearity of the limit, put the finitely
many inputs in one ball, use the same sufficiently large cutoff for all
of them, and apply the linearity equation for that finite cutoff.  The
sum of the explicitly bounded tail errors tends to zero with the displayed
schedule.  Equality in the separated completion gives the desired
linearity equation.  The finite bound \eqref{eq:chain-bound} similarly
passes to the limit, and \eqref{eq:chain-tail} gives
\eqref{eq:chain-limit-tail}.

For clarity, the limiting identity can be checked without invoking an
abstract theorem about an operator space.  For every fixed $x$,
\begin{align*}
 \|(D\mathcal H+\mathcal H D-I)x\|
 &\le \|D(\mathcal H-h_N)x\|
       +\|(\mathcal H-h_N)Dx\|+\|R^Nx\|\\
 &\le \left(1+\frac{2B_DH}{1-\rho}\right)\rho^N\|x\|.
\end{align*}
The last expression is below any prescribed dyadic error by a PR
schedule obtained from \eqref{eq:rational-geometric-block} and a rational
upper bound for its constant factor.  Therefore its represented value is
zero.  This proves \eqref{eq:chain-repaired}.
\end{proof}

The proof produces a single output name by diagonalization.  It does not
assume that every pointwise Cauchy family of operators has a uniform
computable limit.  The uniform majorant \eqref{eq:chain-tail} is the
specific extra information that permits this construction.

\subsection{Only compute the boundary observation that is needed}

A cycle is a vector $x$ satisfying $Dx=0$; a boundary is a vector of
the form $Dy$.  Equation \eqref{eq:chain-repaired} supplies the boundary
witness $y=\mathcal Hx$ for every represented cycle.  No test for the
equality $Dx=0$ is required: the conclusion is conditional on that
equation, while the output recipe is defined for every input.

Often one needs only a boundary approximation.  The exact finite identity
already gives
\begin{equation}
 x-Dh_Nx=h_NDx+R^Nx,
 \qquad
 \|x-Dh_Nx\|
 \le\frac{H}{1-\rho}\|Dx\|+\rho^N\|x\|.
 \label{eq:approximate-cycle}
\end{equation}
This also controls an approximate cycle.  In particular, for $Dx=0$ and
$\|x\|\le2^b$, the shorter cutoff
\begin{equation}
                 N_{\rm boundary}(n,b)=2^s(n+b)
                 \label{eq:boundary-cutoff}
\end{equation}
guarantees $\|x-Dh_Nx\|\le2^{-n}$.  The convention $R^0=I$ makes
the case $n=b=0$ valid as well.  This contract concerns the exact finite
value $h_Nx$; any numerical approximation to it must pay its own input
and evaluation errors through the ledger.

The smaller cutoff does not claim that $h_Nx$ itself is already close to
the full homotopy value.  It is sufficient for the requested boundary
observation, whose error is the residual $R^Nx$.  This is a useful example
of why the observation to be exported should determine the computation.

\subsection{Stable neighborhoods of a certified differential}

Let $D'=D+\Delta$ have the same degree, and suppose the exact structural
equation $(D')^2=0$ is supplied.  Suppose also that
\[
           \|I-Dh-hD\|\le\rho_0,\qquad
           \|h\|\le H,\qquad \|\Delta\|\le\delta.
\]
The new defect is
\begin{equation}
 R'=I-D'h-hD'=R-\Delta h-h\Delta,
 \qquad \|R'\|\le\rho_0+2H\delta.
 \label{eq:chain-perturbation}
\end{equation}
Thus any supplied inequality
$\rho_0+2H\delta\le1-2^{-s}$ allows
Theorem~\ref{thm:chain-repair} to construct a corrected contraction
for $D'$.  The old candidate may itself be only approximate.  If it is
an exact contraction, take $\rho_0=0$.

For a family with a certified bound
$\|D_t-D_u\|\le L|t-u|$, this gives an explicit neighborhood of each
certified parameter: any radius $r$ satisfying
$\rho_0+2HLr\le1-2^{-s}$ works.  It does not provide a contraction
on an entire parameter path unless that path is covered by appropriate
certificates.  Also, a small value of $(D')^2$ cannot replace the equation
$(D')^2=0$: the commutation identity in
Lemma~\ref{lem:finite-chain-correction} uses that equation exactly.

\subsection{A categorical application through a mapping cone}

Let $(C,d_C)$ and $(E,d_E)$ be presented normed chain complexes, and let
$F:C\longrightarrow E$ be a represented bounded chain map, meaning
$d_EF=Fd_C$.  Its mapping cone is the graded direct sum
\[
 \operatorname{Cone}(F)_i=E_i\oplus C_{i-1},\qquad
 D_F=\begin{pmatrix}d_E&F\\0&-d_C\end{pmatrix}.
\]
Use the maximum norm on each direct sum.  The square of this block
operator is zero: its diagonal blocks are $d_E^2,d_C^2$, and its upper
right block is $d_EF-Fd_C$.  Hence this is again a valid differential.
Canonical inclusions and projections have norm at most one.

\begin{theorem}[A defect certificate gives a chain-homotopy inverse]
\label{thm:cone-homotopy-inverse}
Suppose a degree-$+1$ candidate $h$ on $\operatorname{Cone}(F)$ has
certified bounds $\|h\|\le2^c$ and
$\|I-D_Fh-hD_F\|\le1-2^{-s}$.
There is an explicitly represented bounded chain map
$G:E\longrightarrow C$ and explicit bounded homotopies witnessing
$FG\simeq I_E$ and $GF\simeq I_C$.
Their recipes and PR observation schedules are obtained from the one
correction construction above.
\end{theorem}
\begin{proof}
Repair $h$ and write the resulting degree-$+1$ contraction in blocks as
\[
                    \mathcal H=
                    \begin{pmatrix}A&B\\G&J\end{pmatrix}.
\]
Here $A$ and $J$ have degree $+1$, $G$ has degree zero, and $B$ has
degree $+2$, with the shifts determined by the cone.  The lower-left,
upper-left, and lower-right blocks of
$D_F\mathcal H+\mathcal H D_F=I$ are respectively
\begin{align*}
             -d_CG+Gd_E&=0,\\
               d_EA+Ad_E+FG&=I_E,\\
               GF-d_CJ-Jd_C&=I_C.
\end{align*}
The first equation says that $G$ is a chain map.  The second supplies
a homotopy from $I_E$ to $FG$; the third supplies a homotopy from $GF$
to $I_C$.  All three block maps have bounds at most
$2^{s+c}$, since the inclusions and projections have bound one.
Composing their fixed recipes with the output name of $\mathcal H$
gives the stated observation schedules.

At a finite cutoff the same block calculation has an error block of
$R^N$ on the right.  Each block has norm at most $\rho^N$.  Thus the
construction can also export finite approximate inverse and homotopy
equations without first materializing their limiting maps.
\end{proof}

For completeness, the categorical relation used here has an entirely
explicit presentation.  A prescribed homotopy from chain maps $f$ to $g$
is a degree-$+1$ map $K$ with $f-g=dK+Kd$.  The zero map witnesses
reflexivity, $-K$ reverses a homotopy, and addition composes two
homotopies.  If $u$ and $v$ are chain maps, then
\[
 u(f-g)=d(uK)+(uK)d,\qquad
 (f-g)v=d(Kv)+(Kv)d.
\]
Thus homotopy is a congruence for composition, and the theorem produces
an inverse arrow in the corresponding homotopy category.  This is a
scheme for fixed presented complexes and maps, not an internally coded
category containing every Banach space or every bounded operator.

The certificate is stable under simultaneous perturbations of both
complexes and of the map.  Write their changes as
$\Delta_C,\Delta_E,\Delta_F$, with bounds
$\delta_C,\delta_E,\delta_F$.  Assume the new differentials square to
zero and the new map is a chain map.  The cone differential changes by
\[
 \Delta_{\rm cone}=
   \begin{pmatrix}\Delta_E&\Delta_F\\0&-\Delta_C\end{pmatrix},
 \qquad
 \|\Delta_{\rm cone}\|
       \le\max(\delta_E+\delta_F,\delta_C).
\]
Substitution in \eqref{eq:chain-perturbation} gives a single rational
margin test for the perturbed chain-homotopy equivalence.  This preserves
the actual witness and its computational bound, rather than appealing to
an unspecified invariance theorem for homology groups.

\subsection{Continuous logic and the indispensable margin}

There is a direct connection to the continuous-predicate viewpoint
\cite{bbhu2008}.  For fixed $D,h$, the bounded predicate
\[
 x\longmapsto\min\{1,\|x-Dhx-hDx\|\}
                     \quad (\|x\|\le1)
\]
has a continuity modulus obtained from the bounds on $D$ and $h$.
A certified upper bound $\rho<1$ for this predicate on the whole unit
ball is exactly the sort of positive margin used by the correction
theorem.  On a represented finite dimensional compact ball, finite nets
and an explicit modulus can supply such a universal bound from finitely
many observations, including their discretization error.  For rational
matrices with the $\ell^1$ norm, one can instead compute the bound
exactly as the maximum absolute column sum.

An infinite dimensional Banach unit ball is not assumed to have finite
nets.  In that setting the universal bound must come from an operator
estimate or another supplied proof.  Neither this application nor the
existing fixed-formula continuous semantics supplies general compactness,
arbitrary model existence, or a finite net for every bounded metric sort.

The need for a margin is already visible in the two-term complex
\[
             0\longrightarrow\mathbb R
                   \xrightarrow{\;t\;}\mathbb R
                       \longrightarrow0.
\]
For a positive rational $t$, division by $t$ gives a contraction; a
certificate $t\ge2^{-a}$ bounds it by $2^a$.  At $t=0$, the
differential is zero and the contraction equation would say $0=I$.
Moreover, a contraction of norm at most $M$ for positive $t$ must satisfy
$1\le tM$.  Thus no uniformly bounded family of contractions can
extend through $t=0$.  The framework resolves the stable case and
exposes this genuine boundary; a different encoding does not eliminate it.

\subsection{Exact arithmetic and Lean obligations}

For rational finite matrices, all core calculations above are bounded
loops over rational entries, with norm bounds from finite sums and maxima.
For a fixed infinite core presentation, the required linearity, norm, and
differential laws remain explicit hypotheses until their iPRA proofs are
provided.  The finite induction lemmas concern equality or rational
inequality of encoded terms.  The completed-space theorem is assembled
from those finite contracts at each requested observation.  It is not
obtained by applying iPRA induction to a statement quantifying over all
Cauchy names.

One can make the finite invariant especially concrete.  Fix a core vector
$x$ and carry the four states
\[
 u_j=R^jx,\quad v_j=R^jDx,\quad
 \sigma_j=\sum_{i<j}R^ix,\quad
 \tau_j=\sum_{i<j}R^iDx.
\]
They are updated by $u_{j+1}=Ru_j$, $v_{j+1}=Rv_j$ and by adding
$u_j,v_j$ to the two sums.  The invariant includes
$Du_j=v_j$, $D\sigma_j=\tau_j$,
$(I-R)\sigma_j=x-u_j$, and the rational norm bounds on these finite
states.  Its step uses the supplied laws at the current core vectors.
This is a quantifier-free invariant with $x$ as a free parameter;
universal generalization over $x$ occurs after its induction proof.

An exact example uses degrees $0,1,1,2$ on four coordinates and
\[
 D=\begin{pmatrix}
 0&1&0&0\\0&0&0&0\\0&0&0&1\\0&0&0&0
 \end{pmatrix},\qquad
 h=\begin{pmatrix}
 0&0&0&0\\3/4&0&0&0\\1/16&0&0&0\\0&1/8&3/4&0
 \end{pmatrix}.
\]
In the $\ell^1$ norm, $\|h\|=13/16$ and
$R=\tfrac14 I-\tfrac{3}{16}\mathsf E_{32}$ has norm $7/16$,
where $\mathsf E_{32}$ is the matrix with its only nonzero entry in
row~3, column~2.  Here $h^2\ne0$ and $hR\ne Rh$.  Nevertheless the
correction gives the exact contraction
\[
 \mathcal H=\begin{pmatrix}
 0&0&0&0\\1&0&0&0\\1/12&0&0&0\\0&-1/12&1&0
 \end{pmatrix}.
\]
The margin parameters $s=1,c=0$ are valid without computing this
closed form.

The file \texttt{ObservableAnalysis/Neumann.lean} checks the exact
identities of Lemma~\ref{lem:finite-chain-correction} over a possibly
noncommutative ring.  Its definitions mirror the defect, geometric sum,
and corrected homotopy in \eqref{eq:chain-defect}.
The residual theorem is named
\texttt{correctedContraction\_residual}.
The remaining formalization steps are specific:
\begin{enumerate}
\item Realize core vectors and operators by the selected arithmetic
  presentation, and prove the degree, linearity, and norm contracts.
\item Prove the finite geometric inequality and the dyadic block estimate
  \eqref{eq:rational-geometric-block} with encoded rational arithmetic.
\item Apply the represented completion construction to the displayed
  diagonal name $y_n$, proving its error bound and extensionality.
\item Define the graded cone by finite tagged products, prove its
  square-zero identity, and extract the three block equations above.
\item Supply object-language iPRA proofs for the chosen presentation
  laws and finite induction instances.  The compiled Lean ring theorem
  is a verified algebraic metatheorem, not by itself an iPRA derivation.
\end{enumerate}

The exact-rational example in \texttt{code/check\_chain\_contraction.py}
uses a four-dimensional graded complex with $h^2\ne0$ and $hR\ne Rh$.
It checks the finite residual identities and compares the scheduled
approximants with an independently computed rational inverse.  The
recorded result in \texttt{evidence/chain\_contraction\_check.json}
is an example verification; the general finite identity is proved in the
Lean file and in Lemma~\ref{lem:finite-chain-correction}.
